% =====================================================================
%  2.5 Atividade com SonarQube
% =====================================================================
\subsection{Atividade com SonarQube}\label{sec:sonarqube}

\subsubsection{O que é o SonarQube e para que serve a análise estática}

O \textbf{SonarQube} é uma ferramenta de \textbf{análise estática de código}: ele examina o
código-fonte \textbf{sem executar a aplicação} e aponta problemas, exibindo os resultados em
\textit{dashboards} \citep{sonar2026}. A análise estática serve para encontrar cedo -- quando a
correção é mais barata -- \textbf{\textit{bugs}} (erros que causam comportamento incorreto),
\textbf{vulnerabilidades} (falhas exploráveis por um atacante), \textbf{\textit{security hotspots}}
(trechos sensíveis que precisam de revisão), \textbf{\textit{code smells}} (problemas de manutenção),
além de duplicações e baixa cobertura de testes. Um exemplo clássico é uma senha escrita direto no
código. No DevOps, o SonarQube entra no \textit{pipeline} como um ``portão de qualidade''
(\textit{Quality Gate}): se o código não atende aos critérios, a entrega é bloqueada \citep{muniz2026}.

\subsubsection{Criação do ambiente e acesso ao painel}

O SonarQube foi instalado em uma VM própria (\texttt{sonarqube}, \texttt{192.168.1.6}, 4 GB de
memória, porta 9000 redirecionada para o Windows). O script de provisionamento automatizou a
instalação: Java 21, criação do usuário \texttt{sonar} (o Elasticsearch interno não roda como
\texttt{root}), ajuste do \textit{kernel} exigido pelo Elasticsearch (\texttt{vm.max\_map\_count}),
download do SonarQube para \texttt{/opt/sonarqube}, criação do serviço \texttt{sonar} no
\texttt{systemd}, liberação da porta 9000 no \textit{firewall} e instalação do \textbf{SonarScanner
CLI}. O painel foi acessado pelo navegador em \url{http://127.0.0.1:9000}, com o usuário
\texttt{admin}/\texttt{admin}, cuja senha o SonarQube obriga a trocar no primeiro acesso.

\subsubsection{Criação do projeto e execução do sonar-scanner}

No painel, o projeto foi criado em \textit{Create Project} $\rightarrow$ \textit{Local project}, com
o nome \textit{Easy Notes} e a chave \texttt{easy-notes}; como método de análise foi escolhido
\textit{Locally}, e o SonarQube gerou um \textbf{\textit{token}} -- uma credencial que o
\textit{scanner} usa para enviar os resultados. Para projetos Java, o \textit{scanner} precisa também
das classes compiladas, por isso o projeto foi clonado e compilado com o Maven antes da análise
(Código~\ref{cod:sonar-scanner}).

\begin{terminal}[caption={Compilação do projeto e execução do sonar-scanner},label={cod:sonar-scanner}]
git clone https://github.com/callicoder/spring-boot-mysql-rest-api-tutorial.git easy-notes
cd easy-notes && mvn -B clean package -DskipTests     # gera target/classes
export SONAR_TOKEN=************                       # token gerado no painel (oculto)
sonar-scanner \
  -Dsonar.projectKey=easy-notes \
  -Dsonar.sources=src/main/java \
  -Dsonar.java.binaries=target/classes \
  -Dsonar.host.url=http://localhost:9000 \
  -Dsonar.token=$SONAR_TOKEN
\end{terminal}

Os parâmetros indicam a chave do projeto (\texttt{projectKey}), onde está o código
(\texttt{sources}), as classes compiladas (\texttt{java.binaries}), o endereço do servidor
(\texttt{host.url}) e a credencial (\texttt{token}; em versões antigas, \texttt{sonar.login}). Ao
final, o \textit{scanner} exibe \texttt{ANALYSIS SUCCESSFUL} e o resultado aparece no
\textit{dashboard}.

\begin{atencao}[Token]
  O \textit{token} funciona como uma senha: ele foi guardado em uma variável de ambiente, não aparece
  em arquivos versionados e foi ocultado nos prints.
\end{atencao}

\subsubsection{Resultados da análise}

% CONFIRA: preencha a tabela com os números do SEU dashboard (os mesmos do print).
\begin{table}[H]
  \centering
  \caption{Resultado da análise do projeto Easy Notes no SonarQube}\label{tab:resultado-sonar}
  \small
  \renewcommand{\arraystretch}{1.2}
  \begin{tabular}{l c}
    \toprule
    \textbf{Métrica} & \textbf{Resultado} \\
    \midrule
    Quality Gate & \preencher{Passed / Failed} \\
    \textit{Bugs} / Confiabilidade & \preencher{qtd. e nota} \\
    Vulnerabilidades / Segurança & \preencher{qtd. e nota} \\
    \textit{Security Hotspots} & \preencher{qtd.} \\
    \textit{Code Smells} / Manutenibilidade & \preencher{qtd. e nota} \\
    Duplicações e cobertura & \preencher{\% e \%} \\
    \bottomrule
  \end{tabular}
  \fonte{Elaborado pelo autor (2026), a partir do \textit{dashboard} do SonarQube.}
\end{table}

% CONFIRA: ajuste a lista abaixo com os apontamentos que aparecerem na aba Issues do SEU projeto.
Na aba \textit{Issues}, os apontamentos foram principalmente de manutenção e segurança, como: a
\textbf{injeção de dependência por campo} (\texttt{@Autowired} no atributo do
\texttt{NoteController}, quando o recomendado é injetar pelo construtor); a \textbf{entidade JPA
recebida diretamente na API} (\texttt{@RequestBody Note}), risco de \textit{mass assignment} que se
resolve com um objeto de transferência (DTO); o retorno com tipo genérico curinga
(\texttt{ResponseEntity<?>}); e atributos não \texttt{final} na classe de exceção.

\printtela{sonar-token}{Geração do token de análise (valor ocultado)}%
  {Tela de geração do \textit{token} com o valor \textbf{borrado/tarjado}.}

\printtela{sonar-scanner-exec}{Execução do sonar-scanner}%
  {Final da saída do \texttt{sonar-scanner} com \texttt{ANALYSIS SUCCESSFUL} e
   \texttt{EXECUTION SUCCESS} (sem o \textit{token} visível).}

\printtela{sonar-dashboard}{Dashboard do projeto no SonarQube}%
  {Tela \textit{Overview} do projeto com o Quality Gate e as métricas.}

\printtela{sonar-issues}{Problemas apontados na aba Issues}%
  {Aba \textit{Issues} do projeto listando os apontamentos.}

\subsubsection{Como o SonarQube melhora a qualidade e a segurança antes do deploy}

O SonarQube funciona como um fiscal automático: cada versão do código é analisada de forma objetiva,
e os problemas aparecem logo após o \textit{commit}, quando ainda são baratos de corrigir. As regras
aplicadas a cada linguagem formam o \textbf{Quality Profile}, e as condições para aprovar o código
(por exemplo, nenhum problema grave novo e duplicação baixa) formam o \textbf{Quality Gate}, que no
\textit{pipeline} impede que código reprovado chegue ao \textit{deploy}. Ele não substitui testes nem
revisão humana -- não entende a regra de negócio --, mas complementa o processo de qualidade.
